\section{集群、可靠性与数据}
\label{sec:cluster}

把单个作业的并行策略调好以后，大模型训练的真实效率还受三件事影响：多个作业怎样共享集群（调度），万卡规模下的故障怎样才不会吞掉训练进度（可靠性），
以及 GPU 是不是在等数据（数据管线）。这一层是训练 infra 岗位高级面试题的主要来源。

\subsection{深度学习集群调度}

深度学习作业和传统的批处理作业很不一样。Jeon 等人分析微软 Philly 集群~\cite{jeon2019analysis} 后发现：分布式训练需要 gang scheduling，即所有 worker 必须同时启动，否则先启动的只能空等；
作业对拓扑局部性很敏感，最好放在同一节点或同一交换机下，而这种局部性约束又让排队时间变长；还有相当比例的作业因为用户错误或故障而失败。
针对这些特点，Gandiva~\cite{xiao2018gandiva} 利用 DL 作业的迭代性做时间片轮转和迁移，并根据运行时的观察调整放置。Tiresias~\cite{gu2019tiresias} 在不知道作业时长的前提下，
按「已获得的服务量 = GPU 数 $\times$ 时间」做二维 LAS（least attained service）调度，降低平均完成时间。
Pollux~\cite{qiao2021pollux} 注意到吞吐和统计效率都会随 batch 变化，于是联合调整每个作业的 GPU 数和 batch size，以 goodput 为优化目标。

\begin{qa}{深度学习集群调度和传统批处理调度有什么不同？}
\ans 主要有四点不同。第一是 gang scheduling：分布式作业的所有 worker 必须同时拿到资源。第二是拓扑敏感：放在同一节点或同一交换机下，会直接影响 all-reduce 的时间。
第三是可抢占和弹性：作业可以依靠 checkpoint 暂停、迁移和伸缩。第四是吞吐与资源不成线性关系：卡数加倍不一定吞吐加倍，batch 变大还会改变统计效率（Pollux）。
\wrong 「按 GPU 数量排队就行」。
\end{qa}

\subsection{大规模训练的故障与恢复}

到了万卡规模，故障就是常态。Llama~3 报告，在约 1.6 万张 H100 上的一段 54 天预训练中，作业共中断了 466 次，其中 419 次是非计划中断；
约 78\% 归因于已确认的硬件问题（GPU 和主机组件），GPU 相关问题是最大的一类。靠自动化恢复，有效训练时间仍然超过了 90\%~\cite{dubey2024llama3}。
MegaScale~\cite{jiang2024megascale} 在 12,288 张 GPU 上，把算法、通信重叠、算子和数据管线的优化，与「深度可观测性」（细粒度监控、心跳、自动诊断和慢节点剔除）结合起来，
达到了 55.2\% 的 MFU。Gemini 技术报告~\cite{geminiteam2023gemini} 也讨论了大规模训练中的静默数据损坏（SDC）和快速恢复。

\begin{est}{万卡训练多久坏一次？checkpoint 应该多久存一次？}
设单卡的年故障率是 2\%，1 万张卡每年约坏 200 次，平均约 44 小时一次；再算上网络、主机和软件故障，实际往往是十几小时一次。
70B 模型的完整 checkpoint（BF16 权重加 FP32 主权重和 Adam 状态）约为 $16\times70.6\times10^9\approx1.1$~TB。
Young 公式~\cite{young1974first} 给出最优的保存间隔 $\tau\approx\sqrt{2\delta M}$，其中 $\delta$ 是一次保存的开销，$M$ 是平均故障间隔。
取 $\delta=30$~s、$M=12$~h，$\tau\approx\sqrt{2\times30\times43200}\approx1610$~s，也就是约 27 分钟存一次；每次故障平均损失约 $\tau/2$，再加上重启的时间。
\end{est}

从公式可以看出，工程上的关键是把 $\delta$ 压小。异步 checkpoint 先把分片拷到主机内存，每张卡几 GB，几秒就能完成，训练随即继续，再由后台写入分布式存储；
每个 DP 组只需要保存一份；CheckFreq~\cite{mohan2021checkfreq} 则把快照和持久化组成流水线，并自适应地调整保存频率。
$\delta$ 越小，最优间隔就越短，故障造成的平均损失也越小。

\begin{qa}{checkpoint 间隔怎么定？为什么要异步？}
\ans 间隔要在保存开销和故障损失之间平衡：间隔太短，保存开销占比高；太长，故障时丢的进度多。最小化期望损失，就得到 Young 公式 $\tau\approx\sqrt{2\delta M}$。
异步保存把 $\delta$ 从「写分布式存储的时间」降成「拷贝到主机内存的时间」，既降低了开销，又允许存得更频繁。恢复时间也要考虑：从内存副本（或相邻节点上的冗余副本）恢复，比从远端存储读取快得多。
\wrong 「每 N 步存一次就行」，不考虑故障率和保存开销；或者同步写远端存储，把训练阻塞住。
\end{qa}

\begin{qa}{分布式训练挂住了，你怎么查？}
\ans 挂住几乎总是因为集合通信对不上：某个 rank 少调或多调了一次 collective，或者各 rank 的调用顺序不同。我会这样查。先打开 \texttt{NCCL\_DEBUG=INFO} 和异步错误处理，并设一个合理的超时，
让「挂住」变成「报错」。再用 \texttt{py-spy dump} 看每个 rank 停在哪一行，如果停在不同的 collective 上，就是顺序不一致。常见的根因有：只在 rank 0 执行的分支里调用了 collective；
输入长度不同，导致某个 rank 多跑了一轮；部分参数在某些 rank 上没有梯度。如果只是变慢而不是挂住，就先用 \texttt{nccl-tests} 测 bus bandwidth，与式~\eqref{eq:ring} 的理论值比较，
再检查拓扑（是不是走了 PCIe）和 IB 的配置。
\wrong 「重启一下」或「把超时调大」，却不去定位是哪个 rank 不一致。
\end{qa}

\subsection{性能排查的顺序}

训练或推理变慢时，面试官想听的是\emph{排查的顺序}，以及\emph{能区分不同假设的那个观测}。一个可以复述的顺序是：先用 roofline 算出上界，实测如果已经在上界的 1.2 倍以内，就别再优化 kernel 了；
再看时间线上有没有空洞；然后看单个 kernel 的 DRAM 和 SM 吞吐；最后才调超参，而且每次只动一个变量。表~\ref{tab:debug} 把常见的时间线症状和对应的原因、区分方法列在一起。

\begin{table}[!t]
\centering
\caption{时间线症状、最可能的原因与区分性观测}
\label{tab:debug}
\footnotesize
\begin{tabularx}{\columnwidth}{@{}L{2.2cm} L{2.4cm} Y@{}}
\toprule
时间线上看到的 & 最可能的原因 & 区分它的观测 \\
\midrule
kernel 间大量小空洞 & CPU 发射开销、小 kernel 过多 & 空洞宽度与 kernel 数成正比；CUDA Graph 或 \texttt{torch.compile} 后消失 \\
每步开头一段空白 & dataloader 跟不上 & 加倍 \texttt{num\_workers} 后变短；GPU 空闲时 CPU 满载 \\
反向后一大段 NCCL & 梯度通信未与反向重叠 & bucket 过大或开启了 \texttt{find\_unused\_} \texttt{parameters} \\
周期性长尖峰 & 分配器碎片整理、checkpoint & \texttt{num\_alloc\_retries} 增长 \\
利用率 100\% 但慢 & 访存受限 kernel & DRAM 吞吐近峰值、SM 吞吐低 \\
某卡总是最后到 & 负载不均或慢卡 & 每 rank 步时分布；换卡后是否跟随 \\
\bottomrule
\end{tabularx}
\end{table}

\begin{qa}{MFU 只有 25\%，你怎么查？}
\ans 先算账，再看 profile。算账是指用 $6N$ 加上 attention 项估计模型 FLOPs，并确认 MFU 的分母用的是稠密峰值，不含稀疏。然后按占比从大到小逐项排查：通信有没有暴露（没有和计算重叠的 all-reduce 或 all-gather）；
流水线气泡是不是太大（$m<4p$）；dataloader 有没有造成每步开头的空白；逐元素 kernel 有没有融合；micro-batch 是不是太小，GEMM 吃不满；是不是开了全量重计算（这时 HFU 高而 MFU 低）。
每次只改一个因素，改完重新测。
\wrong 一上来就去手写 kernel，或者只报 HFU。
\end{qa}

\subsection{数据管线}

最后一个容易被忽视的问题是：GPU 会不会在等数据？Mohan 等人~\cite{mohan2021analyzing} 发现，很多训练作业有相当一部分时间耗在「数据停顿」（data stall）上，也就是数据读取和预处理跟不上 GPU。
tf.data~\cite{murray2021tfdata} 用流水化、并行 map 和自动调优来缓解这个问题。LLM 预训练的数据管线以离线预分词、内存映射的二进制分片和按长度打包（packing）为主，在线开销较小；
但多模态（要解码图像和视频）和 RL（要生成 rollout）又把数据这一侧推回了关键路径。

这里还有两个经典系统值得一提。Spark 的 RDD~\cite{zaharia2012resilient} 靠 lineage 重算实现容错，是大规模离线数据处理的基础；参数服务器~\cite{li2014scaling} 则在推荐场景中继续承担稀疏 embedding 的训练。
把这两者和 all-reduce 对照着讲，是 CS149 和 15-849 之间最自然的衔接：它们分别代表了数据并行原语在批处理、稀疏训练和稠密训练中的三种形态。
另一个方向是用学习的方法来优化系统本身，比如用强化学习做设备放置~\cite{mirhoseini2017device}，这对研究岗是加分项。

\begin{keypts}
\begin{itemize}
\item DL 调度的特点：gang、拓扑敏感、可抢占与弹性、吞吐与资源非线性（Gandiva、Tiresias、Pollux）。
\item 万卡训练的故障按小时计；Llama~3 的 54 天里有 419 次非计划中断，约 78\% 是硬件问题。
\item Young 公式 $\tau\approx\sqrt{2\delta M}$；用异步、分片、先写内存再落盘来压小 $\delta$。
\item 排查顺序：先算 roofline 上界，再看时间线空洞，然后看单 kernel 吞吐，最后调超参；训练挂住时先查 collective 的调用顺序。
\end{itemize}
\end{keypts}

\begin{research}
\begin{itemize}
\item 静默数据损坏（SDC）在万卡规模下怎样检测？一个可做的设计是周期性地让不同的 rank 重复计算同一个小批次，并比较梯度的哈希；
需要量化它的开销和检出率之间的权衡。
\item 用本文的事件模拟器研究「慢节点」对同步训练的影响：一个 rank 慢 10\% 时，DP、PP、TP 各自的整体吞吐下降多少？
哪种并行方式对尾部延迟最敏感？
\end{itemize}
\end{research}
